% One numbered chain is reused throughout all three sweeps.
% Outline the current edge: yellow if proper, its shared colour if monochromatic.
\makeatletter
\newcommand{\outlineedge}[4]{%
  \ifnum#4=\outlinecurrent\relax
    % Read the actual column colours, including monochromatic edges that Y skips.
    \expandafter\let\expandafter\outlinefirstcolor
      \csname shiftchain@vertexcolor@#2\endcsname
    \let\outlineborder\outlinefirstcolor
    \@for\outlinevertex:=#1\do{%
      \expandafter\ifx\csname shiftchain@vertexcolor@\outlinevertex\endcsname\outlinefirstcolor
      \else\def\outlineborder{cfyel}\fi
    }%
    \draw[draw=\outlineborder,line width=1.6pt,rounded corners=3pt]
      ({#2-0.35},{-#4-0.35}) rectangle ({#3+0.35},{-#4+0.35});
  \fi
  % Trial dots use a pale colour; the outline still reflects the attempted colour.
  \def\outlinedots{}\def\outlineseparator{}%
  \@for\outlinevertex:=#1\do{%
    \edef\outlinedots{\outlinedots\outlineseparator\outlinevertex}%
    \ifnum\outlinevertex=\outlinetrialvertex\relax
      \edef\outlinedots{\outlinedots/\outlinetrialcolour}%
    \fi
    \def\outlineseparator{,}%
  }%
  \expandafter\shiftedge\expandafter{\outlinedots}[$e_{#4}$]
}
\makeatother
\newcommand{\outlinetrial}[2]{%
  \def\outlinetrialvertex{#1}\def\outlinetrialcolour{#2}%
}
% Optional trial styling; bars, labels, red vertices, blue vertices, current edge,
% status, sweep direction (1 = LR, -1 = RL).
\newcommand{\outlinediagram}[8][]{%
  \def\outlinetrialvertex{0}\def\outlinetrialcolour{cfpink}%
  #1%
  \tikzset{outline moved/.style={draw=cfnavy!60!black,line width=1.2pt},
    outline old/.style={dashed,draw=cfgrey,line width=0.6pt}}%
  \def\outlinecurrent{#6}%
  \begin{shiftchain}[
    column sep=5.5mm,row sep=3.5mm,dot radius=2.6pt,
    label font=\scriptsize,bars={#2},bar labels={#3},
    bar style={cfnavy,line width=0.6pt}
  ]{20}
    % Fix the bounds so the chain never jumps between overlays.
    \path[use as bounding box] (-0.8,-11.6) rectangle (20.7,2.5);
    \node at (10.5,2) {#7};
    \ifnum#8=1\relax
      \draw[->,cfnavy,thick] (1,1.1) -- (20,1.1);
    \else
      \draw[<-,cfnavy,thick] (1,1.1) -- (20,1.1);
    \fi
    \shiftvertexcolor{#4}{cfred}
    \shiftvertexcolor{#5}{cfnavy}
    \outlineedge{1,2,3,4}{1}{4}{1}
    \outlineedge{2,3,6,7}{2}{7}{2}
    \outlineedge{4,5,6,7}{4}{7}{3}
    \outlineedge{4,5,6,10}{4}{10}{4}
    \outlineedge{7,8,9,10}{7}{10}{5}
    \outlineedge{11,12,13,14}{11}{14}{6}
    \outlineedge{12,13,16,17}{12}{17}{7}
    \outlineedge{14,15,16,17}{14}{17}{8}
    \outlineedge{14,15,16,20}{14}{20}{9}
    \outlineedge{17,18,19,20}{17}{20}{10}
  \end{shiftchain}%
}

\begin{frame}[t]{Algorithm outline}
  \small\setstretch{1.0}
  \textbf{X:} Colour greedily $\longrightarrow$; switch only to avoid a monochromatic edge with no switches.\par
  \uncover<28->{\textbf{Y:} Sweep $\longleftarrow$; move the first switch of each $(3,4)$-edge just before its third vertex.}\par
  \uncover<40->{\textbf{Z:} Sweep $\longrightarrow$; add a switch just after the third vertex of each remaining monochromatic edge.}\par
  \smallskip
  \centering
   \centering
 \only<1>{
    \begin{shiftchain}[
      column sep=6mm,
      row sep=4mm,
      dot radius=1.8pt,
      dot color=black,
      label font=\scriptsize,
      % bars={3,6/dashed,8},
      % Labels sit beneath the bars, including individually dashed bars.
      % bar labels={3/$x_1$,6/$y_2$,8/$x_3$},
      bar style={cfnavy,line width=0.6pt}
    ]{10}
      % Vertex colours apply to these columns in subsequent rows.
      % \shiftvertexcolor{1,2,3,7,8}{cfred}
      % \shiftvertexcolor{4,5,6,9,10}{cfblue}
      \shiftedge{1,2,3,4}
      \shiftedge{2,3,6,7}
      \shiftedge{4,5,6,7}
      \shiftedge{4,5,6,10}
      \shiftedge{7,8,9,10}
    \end{shiftchain}
}
 \only<2>{
    \begin{shiftchain}[
      column sep=6mm,
      row sep=4mm,
      dot radius=1.8pt,
      dot color=black,
      label font=\scriptsize,
      % bars={3,6/dashed,8},
      % Labels sit beneath the bars, including individually dashed bars.
      % bar labels={3/$x_1$,6/$y_2$,8/$x_3$},
      bar style={cfnavy,line width=0.6pt}
    ]{10}
      % Vertex colours apply to these columns in subsequent rows.
      \shiftvertexcolor{1}{cfred}
      % \shiftvertexcolor{4,5,6,9,10}{cfblue}
      \shiftedge{1,2,3,4}
      \shiftedge{2,3,6,7}
      \shiftedge{4,5,6,7}
      \shiftedge{4,5,6,10}
      \shiftedge{7,8,9,10}
    \end{shiftchain}
}
 \only<3>{
    \begin{shiftchain}[
      column sep=6mm,
      row sep=4mm,
      dot radius=1.8pt,
      dot color=black,
      label font=\scriptsize,
      % bars={3,6/dashed,8},
      % Labels sit beneath the bars, including individually dashed bars.
      % bar labels={3/$x_1$,6/$y_2$,8/$x_3$},
      bar style={cfnavy,line width=0.6pt}
    ]{10}
      % Vertex colours apply to these columns in subsequent rows.
      \shiftvertexcolor{1,2}{cfred}
      % \shiftvertexcolor{4,5,6,9,10}{cfblue}
      \shiftedge{1,2,3,4}
      \shiftedge{2,3,6,7}
      \shiftedge{4,5,6,7}
      \shiftedge{4,5,6,10}
      \shiftedge{7,8,9,10}
    \end{shiftchain}
}
 \only<4>{
    \begin{shiftchain}[
      column sep=6mm,
      row sep=4mm,
      dot radius=1.8pt,
      dot color=black,
      label font=\scriptsize,
      % bars={3,6/dashed,8},
      % Labels sit beneath the bars, including individually dashed bars.
      % bar labels={3/$x_1$,6/$y_2$,8/$x_3$},
      bar style={cfnavy,line width=0.6pt}
    ]{10}
      % Vertex colours apply to these columns in subsequent rows.
      \shiftvertexcolor{1,2,3}{cfred}
      % \shiftvertexcolor{4,5,6,9,10}{cfblue}
      \shiftedge{1,2,3,4}
      \shiftedge{2,3,6,7}
      \shiftedge{4,5,6,7}
      \shiftedge{4,5,6,10}
      \shiftedge{7,8,9,10}
    \end{shiftchain}
}
 \only<5>{
    \begin{shiftchain}[
      column sep=6mm,
      row sep=4mm,
      dot radius=1.8pt,
      dot color=black,
      label font=\scriptsize,
      % bars={3,6/dashed,8},
      % Labels sit beneath the bars, including individually dashed bars.
      % bar labels={3/$x_1$,6/$y_2$,8/$x_3$},
      bar style={cfnavy,line width=0.6pt}
    ]{10}
      % Vertex colours apply to these columns in subsequent rows.
    \shiftvertexcolor{1,2,3}{cfred}
      \shiftvertexcolor{4}{cfpink}
      % \shiftvertexcolor{4,5,6,9,10}{cfblue}
      \shiftedge{1,2,3,4}
      \shiftedge{2,3,6,7}
      \shiftedge{4,5,6,7}
      \shiftedge{4,5,6,10}
      \shiftedge{7,8,9,10}
    \end{shiftchain}
}
 \only<6>{
    \begin{shiftchain}[
      column sep=6mm,
      row sep=4mm,
      dot radius=1.8pt,
      dot color=black,
      label font=\scriptsize,
      bars={3},
      % Labels sit beneath the bars, including individually dashed bars.
      bar labels={3/$x_1$},
    ]{10}
      % Vertex colours apply to these columns in subsequent rows.
    \shiftvertexcolor{1,2,3}{cfred}
      \shiftvertexcolor{4}{cfnavy}
      % \shiftvertexcolor{4,5,6,9,10}{cfblue}
      \shiftedge{1,2,3,4}
      \shiftedge{2,3,6,7}
      \shiftedge{4,5,6,7}
      \shiftedge{4,5,6,10}
      \shiftedge{7,8,9,10}
    \end{shiftchain}
}
 \only<7>{
    \begin{shiftchain}[
      column sep=6mm,
      row sep=4mm,
      dot radius=1.8pt,
      dot color=black,
      label font=\scriptsize,
      bars={3},
      % Labels sit beneath the bars, including individually dashed bars.
      bar labels={3/$x_1$},
    ]{10}
      % Vertex colours apply to these columns in subsequent rows.
    \shiftvertexcolor{1,2,3}{cfred}
      \shiftvertexcolor{4,5,6}{cfnavy}
      % \shiftvertexcolor{4,5,6,9,10}{cfblue}
      \shiftedge{1,2,3,4}
      \shiftedge{2,3,6,7}
      \shiftedge{4,5,6,7}
      \shiftedge{4,5,6,10}
      \shiftedge{7,8,9,10}
    \end{shiftchain}
}
 \only<8>{
    \begin{shiftchain}[
      column sep=6mm,
      row sep=4mm,
      dot radius=1.8pt,
      dot color=black,
      label font=\scriptsize,
      bars={3},
      % Labels sit beneath the bars, including individually dashed bars.
      bar labels={3/$x_1$},
    ]{10}
      % Vertex colours apply to these columns in subsequent rows.
    \shiftvertexcolor{1,2,3}{cfred}
      \shiftvertexcolor{4,5,6}{cfnavy}

      \shiftvertexcolor{7}{cfblue}
      % \shiftvertexcolor{4,5,6,9,10}{cfblue}
      \shiftedge{1,2,3,4}
      \shiftedge{2,3,6,7}
      \shiftedge{4,5,6,7}
      \shiftedge{4,5,6,10}
      \shiftedge{7,8,9,10}
    \end{shiftchain}
}
 \only<9>{
    \begin{shiftchain}[
      column sep=6mm,
      row sep=4mm,
      dot radius=1.8pt,
      dot color=black,
      label font=\scriptsize,
      bars={3,6},
      % Labels sit beneath the bars, including individually dashed bars.
      bar labels={3/$x_1$, 6/$x_2$},
    ]{10}
      % Vertex colours apply to these columns in subsequent rows.
    \shiftvertexcolor{1,2,3,7}{cfred}
      \shiftvertexcolor{4,5,6}{cfnavy}

      % \shiftvertexcolor{7}{cfblue}
      % \shiftvertexcolor{4,5,6,9,10}{cfblue}
      \shiftedge{1,2,3,4}
      \shiftedge{2,3,6,7}
      \shiftedge{4,5,6,7}
      \shiftedge{4,5,6,10}
      \shiftedge{7,8,9,10}
    \end{shiftchain}
}
 \only<10>{
    \begin{shiftchain}[
      column sep=6mm,
      row sep=4mm,
      dot radius=1.8pt,
      dot color=black,
      label font=\scriptsize,
      bars={3,6},
      % Labels sit beneath the bars, including individually dashed bars.
      bar labels={3/$x_1$, 6/$x_2$},
    ]{10}
      % Vertex colours apply to these columns in subsequent rows.
    \shiftvertexcolor{1,2,3,7,8,9}{cfred}
      \shiftvertexcolor{4,5,6}{cfnavy}

      % \shiftvertexcolor{7}{cfblue}
      % \shiftvertexcolor{4,5,6,9,10}{cfblue}
      \shiftedge{1,2,3,4}
      \shiftedge{2,3,6,7}
      \shiftedge{4,5,6,7}
      \shiftedge{4,5,6,10}
      \shiftedge{7,8,9,10}
    \end{shiftchain}
}
 \only<11>{
    \begin{shiftchain}[
      column sep=6mm,
      row sep=4mm,
      dot radius=1.8pt,
      dot color=black,
      label font=\scriptsize,
      bars={3,6,9},
      % Labels sit beneath the bars, including individually dashed bars.
      bar labels={3/$x_1$, 6/$x_2$, 9/$x_3$},
    ]{10}
      % Vertex colours apply to these columns in subsequent rows.
    \shiftvertexcolor{1,2,3,7,8,9}{cfred}
      \shiftvertexcolor{4,5,6,10}{cfnavy}

      % \shiftvertexcolor{7}{cfblue}
      % \shiftvertexcolor{4,5,6,9,10}{cfblue}
      \shiftedge{1,2,3,4}
      \shiftedge{2,3,6,7}
      \shiftedge{4,5,6,7}
      \shiftedge{4,5,6,10}
      \shiftedge{7,8,9,10}
    \end{shiftchain}
}

   \only<12>{
    \begin{shiftchain}[
      column sep=6mm,
      row sep=4mm,
      dot radius=1.8pt,
      dot color=black,
      label font=\scriptsize,
      bars={3,6,9},
      % Labels sit beneath the bars, including individually dashed bars.
      bar labels={3/$x_1$, 6/$x_2$, 9/$y_3$},
    ]{10}
      % Vertex colours apply to these columns in subsequent rows.
    \shiftvertexcolor{1,2,3,7,8,9}{cfred}
      \shiftvertexcolor{4,5,6,10}{cfnavy}

      % \shiftvertexcolor{7}{cfblue}
      % \shiftvertexcolor{4,5,6,9,10}{cfblue}
      \shiftedge{1,2,3,4}
      \shiftedge{2,3,6,7}
      \shiftedge{4,5,6,7}
      \shiftedge{4,5,6,10}
      \shiftedge{7,8,9,10}
    \end{shiftchain}
}
 \only<13>{
    \begin{shiftchain}[
      column sep=6mm,
      row sep=4mm,
      dot radius=1.8pt,
      dot color=black,
      label font=\scriptsize,
      bars={3,5,6/dashed,9},
      % Labels sit beneath the bars, including individually dashed bars.
      bar labels={3/$x_1$, 5/$y_2$, 6/$x_2$, 9/$y_3$},
    ]{10}
      % Vertex colours apply to these columns in subsequent rows.
    \shiftvertexcolor{1,2,3,7,8,9}{cfred}
      \shiftvertexcolor{4,5,10}{cfnavy}

      \shiftvertexcolor{6}{cfpink}
      % \shiftvertexcolor{4,5,6,9,10}{cfblue}
      \shiftedge{1,2,3,4}
      \shiftedge{2,3,6,7}
      \shiftedge{4,5,6,7}
      \shiftedge{4,5,6,10}
      \shiftedge{7,8,9,10}
    \end{shiftchain}
}


 \only<14>{
    \begin{shiftchain}[
      column sep=6mm,
      row sep=4mm,
      dot radius=1.8pt,
      dot color=black,
      label font=\scriptsize,
      bars={3,5,6/dashed,9},
      % Labels sit beneath the bars, including individually dashed bars.
      bar labels={3/$x_1$, 5/$y_2$, 6/$x_2$, 9/$y_3$},
    ]{10}
      % Vertex colours apply to these columns in subsequent rows.
    \shiftvertexcolor{1,2,3,6,7,8,9}{cfred}
      \shiftvertexcolor{4,5,10}{cfnavy}

      % \shiftvertexcolor{4,5,6,9,10}{cfblue}
      \shiftedge{1,2,3,4}
      \shiftedge{2,3,6,7}
      \shiftedge{4,5,6,7}
      \shiftedge{4,5,6,10}
      \shiftedge{7,8,9,10}
    \end{shiftchain}
}




  
\only<15>{
    \begin{shiftchain}[
      column sep=6mm,
      row sep=4mm,
      dot radius=1.8pt,
      dot color=black,
      label font=\scriptsize,
      bars={3,5,9},
      % Labels sit beneath the bars, including individually dashed bars.
      bar labels={3/$x_1$, 5/$y_2$, 9/$y_3$},
    ]{10}
      % Vertex colours apply to these columns in subsequent rows.
    \shiftvertexcolor{1,2,3,6,7,8,9}{cfred}
      \shiftvertexcolor{4,5,10}{cfnavy}

      % \shiftvertexcolor{4,5,6,9,10}{cfblue}
      \shiftedge{1,2,3,4}
      \shiftedge{2,3,6,7}
      \shiftedge{4,5,6,7}
      \shiftedge{4,5,6,10}
      \shiftedge{7,8,9,10}
    \end{shiftchain}
}
\only<16>{
    \begin{shiftchain}[
      column sep=6mm,
      row sep=4mm,
      dot radius=1.8pt,
      dot color=black,
      label font=\scriptsize,
      bars={3,5,6,9},
      % Labels sit beneath the bars, including individually dashed bars.
      bar labels={3/$x_1$, 5/$y_2$,6/$y_3$, 9/$y_3$},
    ]{10}
      % Vertex colours apply to these columns in subsequent rows.
    \shiftvertexcolor{1,2,3,6,10}{cfred}
      \shiftvertexcolor{4,5}{cfnavy}

      \shiftvertexcolor{7,8,9}{cfblue}

      % \shiftvertexcolor{4,5,6,9,10}{cfblue}
      \shiftedge{1,2,3,4}
      \shiftedge{2,3,6,7}
      \shiftedge{4,5,6,7}
      \shiftedge{4,5,6,10}
      \shiftedge{7,8,9,10}
    \end{shiftchain}
}
\only<17>{
    \begin{shiftchain}[
      column sep=6mm,
      row sep=4mm,
      dot radius=1.8pt,
      dot color=black,
      label font=\scriptsize,
      bars={3,5,6,9},
      % Labels sit beneath the bars, including individually dashed bars.
      bar labels={3/$x_1$, 5/$y_2$,6/$y_3$, 9/$y_3$},
    ]{10}
      % Vertex colours apply to these columns in subsequent rows.
    \shiftvertexcolor{1,2,3,6,10}{cfred}
      \shiftvertexcolor{4,5,7,8,9}{cfnavy}

      % \shiftvertexcolor{4,5,6,9,10}{cfblue}
      \shiftedge{1,2,3,4}
      \shiftedge{2,3,6,7}
      \shiftedge{4,5,6,7}
      \shiftedge{4,5,6,10}
      \shiftedge{7,8,9,10}
    \end{shiftchain}
}

\end{frame}
